\begin{document}


\lhead[ \leftmark ]{Informatik und Gesellschaft}
\rhead[Informatik]{Cyril Blum, \today, Winterthur}

\part*{Vertiefungsthema}

Im Rahmen dieses Moduls vertiefen Sie sich in ein bestimmtes Thema an der Schnittstelle von Informatik und Gesellschaft. Am Ende des Moduls schreiben Sie einen individuellen Aufsatz, in welchem Sie Ihre Meinung zum Thema als pointierten \textit{Essay} formulieren.

\section{Themenübersicht}

Bilden Sie sich ein inhaltliches Verständnis sowie eine eigene Meinung zu einer der folgenden Aussagen:

\begin{itemize}
	\item \enquote{Kurzvideo-Plattformen wie TikTok fördern die Oberflächlichkeit und verkürzen die Aufmerksamkeitsspanne junger Menschen.}
	
	\item \enquote{Telegram ist ein Leuchtturm der Meinungsfreiheit – und ein trojanisches Pferd für Kriminalität und Desinformation.}
	\item \enquote{Temu manipuliert Sie gezielt und auf vielfältige Weise, damit Sie mehr Dinge einkaufen.}
	\item \enquote{Plattformen wie Telegram und TikTok sollten stärker reguliert werden, um Menschen in demokratischen Staaten vor Überwachung und Falschinformationen zu schützen.}
	\item \enquote{Die Generation Z ist durch Smartphones und soziale Netzwerke weniger kreativ und anfälliger auf psychische Probleme.}
	\item \enquote{Das Internet befeuert den Klimawandel und verschlimmert globale Umweltprobleme.}
	\item \enquote{Firmen sollten im öffentlichen Interesse gezwungen werden, Algorithmen für Allgemeinheit transparent zu machen.}
\end{itemize}

Zu jeder dieser Aussagen dürfen Sie wahlweise die Befürworter- oder auch die Gegenposition einnehmen. Weitere Themen dürfen Sie in gemeinsamer Absprache vorschlagen. 

\section{Vorbereitungszeit und Essay}

In den ersten beiden Doppellektionen erhalten Sie Zeit, um sich in das Thema einzulesen und Argumente zu sammeln. Verschiedene Artikel dazu finden Sie auf Moodle. Selbstverständlich dürfen Sie zur Vorbereitung auch weitere Quellen zum Thema konsultieren. In der dritten Doppellektion verfassen Sie einen handgeschriebenen Essay zum Thema, den Sie wahlweise am Ende der Doppellektion zur Bewertung abgeben können oder nicht. Dieser wird nach folgenden Aspekten benotet:

\begin{enumerate}
\item \textbf{Formales} (1 Punkt): 
\begin{enumerate}
\item Umfang: mindestens 1, maximal 2 Seiten (handgeschrieben, auf A4-Papier)
\item Sprachliche Korrektheit
\item Eleganter, \enquote{süffiger} und leserfreundlicher Schreib-Stil
\end{enumerate}
\item \textbf{Inhaltliche Richtigkeit} (2 Punkte): Die präsentierten Inhalte stimmen und betreffen die zentrale Fragestellung, bzw. Aussage
\item \textbf{Inhaltliche Vollständigkeit} (2 Punkte): Die wichtigsten Argumente der eigenen Sichtweise werden alle vorgebracht und die wichtigsten Gegenargumente werden ebenfalls diskutiert. Mindestens drei wichtige und klar unterscheidbare Argumente sowie drei unterschiedliche Gegenargumente sollten behandelt werden.
\end{enumerate}

An das Essay mitnehmen dürfen Sie die ausgedruckten und angestrichenen Artikel aber keine weiteren Unterlagen oder Notizen.

Die Arbeit wird mit einer Note von 1-6 bewertet und zählt wie eine zusätzliche, halbe schriftliche Prüfung. Sofern Sie die Arbeit abgeben, erklären Sie sich einverstanden damit, dass diese benotet wird und in Ihre Endnote einfliesst.
\end{document}